\documentclass[11pt,a4paper]{report}
\usepackage[T1]{fontenc}
\usepackage[utf8]{inputenc}
\usepackage{geometry}
\geometry{margin=2.5cm}
\usepackage{xcolor}
\usepackage{amsmath}
\usepackage{hyperref}

% ---- latex_handler theme support ------------------------------------------
% For non-default themes latex_handler injects \def\lhtheme{dark} through
% latexmk's -pretex option. Declare a light fallback and branch on the macro.
\providecommand{\lhtheme}{light}
\def\lhdarkvalue{dark}
\definecolor{lhlightaccent}{HTML}{1F4E79}
\definecolor{lhdarkaccent}{HTML}{E8A33D}
\colorlet{lhaccent}{lhlightaccent}
\def\lhthemevalue{Light}
\ifx\lhtheme\lhdarkvalue
  \colorlet{lhaccent}{lhdarkaccent}
  \def\lhthemevalue{Dark}
\fi
\providecommand{\lhbar}{}
\providecommand{\lhfootertext}{}
\InputIfFileExists{lh-branding}{}{}
% ---------------------------------------------------------------------------

\title{An Example Thesis\\[0.4em]\large Built from \texttt{latex-templates}}
\author{Your Name}
\date{\today}

\begin{document}

\maketitle
\lhbar

\begin{abstract}
This is the flagship thesis/report variation of the \texttt{latex-templates}
template. It demonstrates document discovery by \texttt{latex\_handler}, the
light/dark theme mechanism, a project-local asset, and template-shared
branding. Replace the placeholder text with your own work.
\end{abstract}

\tableofcontents
\clearpage

\chapter{Introduction}
This document is compiled from \texttt{src/main.tex}, which matches the default
discovery pattern \texttt{main.tex}. The other variations in this template use
the \texttt{*\_main.tex} pattern.

\input{demo-note}

\chapter{Background}
Research templates benefit from a consistent layout. This template ships its
own \texttt{latex\_handler.toml} so every generated project starts with the
same directory structure, discovery rules, engine, and enabled themes.

\chapter{Method}
The method can be described with ordinary LaTeX. Inline mathematics such as
$E = mc^2$ works without extra setup, and displayed mathematics is available
through \texttt{amsmath}:
\begin{equation}
  \label{eq:example}
  f(x) = \int_0^x t^2 \, \mathrm{d}t = \frac{x^3}{3}.
\end{equation}

\chapter{Results}
Table~\ref{tab:example} shows a small result table. Reference
Equation~\ref{eq:example} from elsewhere in the document with
\texttt{\textbackslash ref}.

\begin{table}[htbp]
  \centering
  \begin{tabular}{lrr}
    \hline
    Setting & Light & Dark \\
    \hline
    Accent & Blue & Orange \\
    Output stem & \texttt{main} & \texttt{main-dark} \\
    \hline
  \end{tabular}
  \caption{Example results.}
  \label{tab:example}
\end{table}

\chapter{Conclusion}
The active theme for this build is \lhthemevalue. Compile with
\texttt{--theme all} to produce both \texttt{main.pdf} and
\texttt{main-dark.pdf}.

\lhbar
\noindent\small\lhfootertext

\end{document}
